%% sections/04_limitations.tex
%% Module 4: Methodology Limitations and Future Work

\section{Methodology Limitations and Future Work}
\label{sec:limitations}

Quantitative models are only as reliable as the assumptions baked into them. This section
documents the specific limitations of the framework presented in this paper. We are
explicit about what we know, what we are assuming, and where a future study could
improve on this work. Readers who are primarily interested in the conclusions should
treat the scenario outputs as order-of-magnitude reasoning rather than precise forecasts.

\subsection{Private Valuation Proxies and the Float Drain Estimate}

The most significant limitation of this study is the reliance on secondary market
transaction prices and journalistic valuation estimates for the three principal IPO
candidates. SpaceX targets a valuation of \$1.75T--\$2.3T, Anthropic has been
reported at approximately \$1T, and OpenAI at \$850B--\$1.1T
\citep{bloomberg_spacex_2026,wsj_anthropic_ipo,ft_openai_ipo}.

These figures are not equivalent to public market valuations. They are subject to
three structural distortions:

\begin{enumerate}
  \item \textbf{Illiquidity premia:} Shares trading on Forge Global, Hiive, or CartaX
        command a 15--40\% premium relative to theoretical fundamental value, driven
        purely by scarcity. This premium compresses at IPO when supply becomes
        unrestricted \citep{bessembinder2020}.
  \item \textbf{Selection bias:} Only buyers with a positive view on a private company
        participate in secondary markets. Sellers are typically insiders with informational
        advantages. The resulting prices reflect the most optimistic end of the valuation
        distribution, not the consensus.
  \item \textbf{Instrument mismatch:} The most recent private rounds for OpenAI
        (\$40B raised at a \$340B pre-money in March 2025) and Anthropic (\$7.5B at a
        \$61.5B pre-money in November 2024) are structured as preferred equity with
        liquidation preferences. Common equity at IPO is structurally junior and will
        trade at a discount to the preferred-share implied valuation.
\end{enumerate}

A 50\% discount to the stated preferred-round valuations would reduce the aggregate
float supply drain estimate to approximately \$105 billion. That figure is still
large enough to qualify as the most significant synchronised IPO liquidity extraction
in US market history. However, it produces a materially different concentration
amplifier and a correspondingly lower peak PFE across all three scenarios.

\textbf{Practical implication:} Investors and risk managers should maintain parallel
scenario estimates at 50\%, 75\%, and 100\% of the stated float drain. The structural
argument (concentrated supply shock into a high-concentration index) holds at all three
levels. The severity of the tail risk does not.

\subsection{The Flat FCF Baseline and Revenue Uncertainty}

The Monetisation Runway calculation in Equation~\ref{eq:runway} assumes that aggregate
hyperscaler free cash flow remains flat at its Q1~2026 annualised level (\$127B) while
Capex continues at the trailing 32.4\% CAGR. We do not forecast hyperscaler revenue.
The flat FCF assumption is a stress scenario, not a central estimate.

This assumption is violated in several plausible ways:

\begin{itemize}
  \item \textbf{Revenue upside:} AI application revenue materialises faster than the
        Capex curve would imply. Specifically, if enterprise AI software adoption
        accelerates (driven by agentic AI deployment in 2026--2027), operating
        leverage could expand FCF margins rapidly. A 20\% FCF growth rate would
        extend the implied runway by approximately 4 additional quarters.
  \item \textbf{Capex downside:} Hyperscalers accelerate Capex further in response to
        competitive pressure without corresponding revenue growth. This is the scenario
        most consistent with our thesis. Capex-to-FCF ratios above $3\times$
        historically trigger investor-imposed capital discipline through shareholder
        activism or covenant restrictions.
  \item \textbf{Non-linear discontinuity:} A Capex pause triggered by a major GPU
        supply disruption, a regulatory intervention (e.g., export controls on advanced
        chip exports), or a voluntary strategic review by one of the four hyperscalers
        would compress the aggregate Capex trajectory in a way our linear CAGR model
        cannot represent. This would be a bullish discontinuity for our
        FCF-recovery thesis.
\end{itemize}

\subsection{Correlation Matrix Stability and the 2022 Calibration Window}

The \texttt{ScenarioStressTester} calibrates its correlation matrix using the 2022
drawdown window (January through December 2022). This is the most recent analogous
stress period available in our dataset. However, a 2026 liquidity shock driven by
IPO supply would differ structurally from the 2022 rate-driven drawdown in ways that
the historical correlation matrix may not capture.

Three specific structural differences warrant caution:

\begin{enumerate}
  \item \textbf{Cause of drawdown:} The 2022 drawdown was driven by Federal Reserve
        rate hikes affecting the discount rate applied to long-duration growth assets.
        A 2026 IPO-supply shock operates through a different mechanism (liquidity
        absorption and forced selling), which historically produces faster initial
        velocity but a shorter duration.
  \item \textbf{Fixed income correlation:} In 2022, equities and bonds fell together
        (positive correlation) because the rate shock was the common factor. In a
        liquidity-driven shock, the correlation between equities and Treasuries is
        more likely to revert to negative (flight-to-quality), altering the effective
        hedging properties of standard 60/40 portfolios.
  \item \textbf{Compute derivative feedback:} If a functioning GPU futures market
        exists by 2026, delta-hedging activity by compute option market-makers would
        introduce a new mechanical correlation between NDX realised volatility and GPU
        spot prices. This correlation does not exist in our 2022 calibration data. It
        represents an entirely unmodelled risk factor that could either amplify or
        dampen the scenario outcomes depending on the net hedge direction.
\end{enumerate}

\textbf{Practical implication:} The scenarios presented in Section~\ref{sec:impact}
should be re-calibrated with an updated correlation matrix as Q3 2026 approaches
and as compute derivative market pricing becomes observable. The quantitative
infrastructure built in this study supports that re-calibration with no
structural changes to the model.

\subsection{Future Work}

\begin{enumerate}
  \item \textbf{Real-time index weight data:} Replace interpolated concentration
        estimates with actual daily constituent weight data from Bloomberg or Refinitiv.
        The current HHI estimates are based on quarterly published disclosures, which
        understates intra-quarter concentration changes during fast-moving market events.
  \item \textbf{Credit spread propagation:} Extend the VaR/PFE framework to
        investment-grade and high-yield credit spreads, modelling the second-order
        contagion pathway from equity drawdown to corporate bond market liquidity.
        Hyperscalers with investment-grade ratings and large debt programmes represent
        a meaningful channel for systemic risk propagation.
  \item \textbf{OCPI integration:} Incorporate the Ornn Compute Price Index spot pricing
        (or CoreWeave equivalent) as a real-time proxy for AI infrastructure
        monetisation velocity. A sustained decline in GPU rental rates would be an
        early-warning signal for Capex deceleration.
  \item \textbf{Multi-asset PFE:} Generalise the PFE calculation from equity-only to a
        cross-asset portfolio including TIPS, VIX, and sector ETF exposures. This would
        allow the framework to produce hedging-strategy-level output rather than index-level
        stress estimates.
  \item \textbf{Bayesian scenario updating:} Replace the fixed Fizzle, Systemic Short,
        and Dotcom Pop scenario probabilities with a Bayesian framework. As IPO filing
        dates, S-1 prospectus valuations, and macroeconomic data arrive through 2026,
        the posterior probability mass over scenarios can be updated dynamically.
  \item \textbf{Agent-based simulation:} Model the interaction between index passive
        vehicles, leveraged ETF holders, and macro hedge fund short positions during a
        stress event. The linear PFE model currently does not represent the
        feedback loop between forced selling and price impact. An agent-based
        simulation would quantify this non-linearity.
\end{enumerate}
